\clearpage
\section{Overview}

@ARCH_NAME@ defines a bounded, 16-bit-word-oriented CISC instruction set
architecture. It includes selected register-memory forms, compact
register-register forms, explicit instruction lengths, and a fixed maximum
instruction size.

The base architecture combines simplified effective addressing with
page-granular segmented address pre-translation, optional page-table
translation, and ordinary memory-mapped device access.

\subsection{Design Goals}
@ARCH_NAME@ is intended for high-performance general-purpose systems rather
than as a minimal embedded instruction set. The architecture keeps the visible
programming model scalar and bounded, while giving implementations explicit
repeated-operation structures that may be executed by microarchitectural
streaming or internal SIMD machinery.

\begin{itemize}
\item preserve selected CISC code-density advantages without unbounded instruction forms
\item make instruction boundaries and maximum instruction size explicit
\item support register-memory operations and compact hot-path register-register operations
\item expose scalar architectural state while allowing wider internal execution resources
\item accelerate fast, regular loops through REP, REPcc, REPG, and update-eligible effective addresses
\item let common loop kernels scale with implementation width without adding programmer-visible vector state or width-specific opcode families
\item avoid making vector width or internal streaming resources part of the architectural ABI
\item leave specialized cryptographic, image, matrix, tensor, and accelerator workloads to extensions or devices
\end{itemize}

\subsection{Architectural Profile}
\begin{itemize}
\item 16-bit instruction words
\item explicit instruction length encoded in word 0
\item maximum instruction length of eight words
\item at most one prefix word containing two prefix bytes
\item register-register and register-memory instruction forms
\item bounded compact and extended effective-address forms
\item page-granular segment pre-translation before optional paging
\item memory-mapped device access through the normal memory map
\item no separate I/O port address space
\item no direct physical-memory load/store instructions
\end{itemize}

\subsection{Instruction Stream}
The instruction stream is composed of 16-bit words. Instructions are aligned to
16-bit boundaries. The program counter addresses bytes, but instruction length
is defined in units of 16-bit words.

\[
\mathit{next\_pc} = \mathit{pc} + 2 \times \mathit{instruction\_length\_in\_words}
\]

The maximum instruction length is eight words, or 16 bytes. No instruction may
exceed this bound.

\subsection{Word 0 Format}
Every instruction begins with word 0. The instruction boundary is determined
from word 0 alone.

\Needspace{1.35in}
\begin{center}\vspace{3pt}
\begin{tikzpicture}[x=0.300in,y=0.24in,every node/.style={font=\scriptsize}]
\node[anchor=east] at (-0.35,0.50) {word 0};
\node[anchor=south] at (0.50,1.18) {15};
\node[anchor=south] at (1.50,1.18) {14};
\node[anchor=south] at (4.50,1.18) {11};
\node[anchor=south] at (15.50,1.18) {0};
\draw (0,0) rectangle (1,1);
\node at (0.50,0.50) {P};
\draw (1,0) rectangle (4,1);
\node at (2.50,0.50) {L};
\draw (4,0) rectangle (16,1);
\node at (10.00,0.50) {primary opcode/fields};
\end{tikzpicture}
\manualfigurecaption{Word 0 Format}
\end{center}


\manualfield{P:}{Prefix-present bit. If set, word 1 is interpreted as the prefix word.}
\manualfield{L:}{Total instruction length encoded as length minus one.}

\begingroup\footnotesize
\setlength{\LTpre}{2pt}
\setlength{\LTpost}{2pt}
\begin{longtable}{@{}>{\raggedright\arraybackslash}p{1.2in}>{\raggedright\arraybackslash}p{1.35in}>{\raggedright\arraybackslash}p{1.35in}@{}}
\toprule
\textbf{L Field} & \textbf{Total Words} & \textbf{Total Bytes}\\
\midrule
\endhead
\texttt{000} & 1 & 2\\
\texttt{001} & 2 & 4\\
\texttt{010} & 3 & 6\\
\texttt{011} & 4 & 8\\
\texttt{100} & 5 & 10\\
\texttt{101} & 6 & 12\\
\texttt{110} & 7 & 14\\
\texttt{111} & 8 & 16\\
\bottomrule
\end{longtable}
\manualtablecaption{Instruction Length Encoding}
\endgroup


\[
\mathit{length} = \mathit{word0.L} + 1 \qquad
\mathit{next\_pc} = \mathit{pc} + 2 \times \mathit{length}
\]

\subsection{Prefix Model}
If P is set, word 1 is a prefix word containing two independent 8-bit prefix slots.
The low byte, bits 7..0, is \texttt{prefix0}; the high byte, bits 15..8, is
\texttt{prefix1}. Prefixes do not determine instruction length.

@PREFIX_MODEL_TABLE@

\Needspace{1.35in}
\begin{center}\vspace{3pt}
\begin{tikzpicture}[x=0.300in,y=0.24in,every node/.style={font=\scriptsize}]
\node[anchor=east] at (-0.35,0.50) {word 1};
\node[anchor=south] at (0.50,1.18) {15};
\node[anchor=south] at (7.50,1.18) {8};
\node[anchor=south] at (8.50,1.18) {7};
\node[anchor=south] at (15.50,1.18) {0};
\draw (0,0) rectangle (8,1);
\node at (4.00,0.50) {prefix1};
\draw (8,0) rectangle (16,1);
\node at (12.00,0.50) {prefix0};
\end{tikzpicture}
\manualfigurecaption{Prefix Word Format}
\end{center}

